% --- Archon physics formalization source begin ---
% archon:physics
% archon:covers PhyXMiniProblems/problem_phyx_mini_0062.lean
% archon:source-report reports/phyx_mini/problem_phyx_mini_0062.source.json

\chapter{Physics problem phyx\_mini\_0062}
\label{ch:PhyXMiniProblems_problem_phyx_mini_0062}

\paragraph{Problem source.}
\#\# Physical scenario

The laser beam in Figure is aimed at the center of a rotating hexagonal mirror.

\#\# Question

How long is the streak of laser light as the reflected laser beam sweeps across the wall behind the laser?

\#\# Answer choices

- A: A: 1.4m

- B: B: 5.6m

- C: C: 6.1m

- D: D: 2.4m

\#\# Figure caption (auxiliary; use the image as primary evidence)

The image is a diagram showing the setup for a laser experiment. Here's a description of the content:

- On the left side, there is a hexagonal object with a side labeled "40 cm."
- An arrow within the hexagon indicates a rotational motion.
- To the right of the hexagon, there is a rectangle labeled "Laser," with an arrow pointing from it towards the hexagon.
- The distance between the laser and the hexagon is labeled as "50 cm."
- Further right, a vertical line labeled "Wall" is shown as the final point in the setup.
- The overall distance from the hexagonal object to the wall is labeled as "2.0 m."

The scene depicts a laser aimed at a hexagonal object, calculating distances to the wall.

\#\# Dataset metadata

Geometrical Optics; Spatial Relation Reasoning, Physical Model Grounding Reasoning

\paragraph{Recorded answer/context.}
C: C: 6.1m

\paragraph{Figure/image path.}
/root/proposal\_for\_physic/science-mango/phyx\_mini\_run/phyx\_data/test\_image/62.png

\paragraph{Formalization target.}
create a compiling Lean file with sorry bodies at `PhyXMiniProblems/problem\_phyx\_mini\_0062.lean`. The Lean declarations must preserve the physical quantities, dimensions or dimensional roles, figure labels, governing-law hypotheses, and final relation expressed by this problem.
Use Mathlib/Physlib names found through LeanExplore where available. If a physics API is missing, introduce faithful local abstractions rather than scalar placeholder aliases.

\begin{theorem}[Physics formalization target]
\label{thm:physics:phyx_mini_0062:target}
\lean{PhyXMiniProblems.ProblemPhyXMini0062.problem_phyx_mini_0062}
\uses{def:physics:phyx-mini-0062:phyxminiproblems-problemphyxmini0062-rotatinghexagonalmirrorsetup, def:physics:phyx-mini-0062:phyxminiproblems-problemphyxmini0062-hasdepictedopticallayout, def:physics:phyx-mini-0062:phyxminiproblems-problemphyxmini0062-matchesfigurereadouts, def:physics:phyx-mini-0062:phyxminiproblems-problemphyxmini0062-hasphysicaldimensions, def:physics:phyx-mini-0062:phyxminiproblems-problemphyxmini0062-satisfiesregularhexagontransitiongeometry, def:physics:phyx-mini-0062:phyxminiproblems-problemphyxmini0062-satisfieslawofspecularreflection, def:physics:phyx-mini-0062:phyxminiproblems-problemphyxmini0062-satisfieswallprojectiongeometry, def:physics:phyx-mini-0062:phyxminiproblems-problemphyxmini0062-isnearestdisplayedstreaklength}
The assigned autoformalize agent should translate this physics problem into Lean declarations in the covered file, with theorem and lemma proof bodies written as `by sorry`.
\end{theorem}
\begin{proof}
This is an autoformalization task, not a proof task. Produce statements that can later be proved by the physics prover without weakening the physical model.
\end{proof}

% --- Archon declaration topology begin ---
\section{Declaration topology}
\begin{definition}[Length Quantity]
  \label{def:physics:phyx-mini-0062:phyxminiproblems-problemphyxmini0062-lengthquantity}
  \lean{PhyXMiniProblems.ProblemPhyXMini0062.LengthQuantity}
  A physical length or signed one-dimensional displacement
\end{definition}
\begin{proof}
  This is the declared model definition of Length Quantity.
\end{proof}
\begin{definition}[length Readout]
  \label{def:physics:phyx-mini-0062:phyxminiproblems-problemphyxmini0062-lengthreadout}
  \lean{PhyXMiniProblems.ProblemPhyXMini0062.lengthReadout}
  \uses{def:physics:phyx-mini-0062:phyxminiproblems-problemphyxmini0062-lengthquantity}
  The scalar readout of a dimensionful length in a selected length unit
\end{definition}
\begin{proof}
  This is the declared model definition of length Readout.
\end{proof}
\begin{definition}[length In Meters]
  \label{def:physics:phyx-mini-0062:phyxminiproblems-problemphyxmini0062-lengthinmeters}
  \lean{PhyXMiniProblems.ProblemPhyXMini0062.lengthInMeters}
  \uses{def:physics:phyx-mini-0062:phyxminiproblems-problemphyxmini0062-lengthquantity, def:physics:phyx-mini-0062:phyxminiproblems-problemphyxmini0062-lengthreadout}
  The scalar SI readout of a dimensionful length, in meters
\end{definition}
\begin{proof}
  This is the declared model definition of length In Meters.
\end{proof}
\begin{definition}[length In Centimeters]
  \label{def:physics:phyx-mini-0062:phyxminiproblems-problemphyxmini0062-lengthincentimeters}
  \lean{PhyXMiniProblems.ProblemPhyXMini0062.lengthInCentimeters}
  \uses{def:physics:phyx-mini-0062:phyxminiproblems-problemphyxmini0062-lengthquantity, def:physics:phyx-mini-0062:phyxminiproblems-problemphyxmini0062-lengthreadout}
  The scalar centimeter readout of a dimensionful length
\end{definition}
\begin{proof}
  This is the declared model definition of length In Centimeters.
\end{proof}
\begin{definition}[Mirror Shape]
  \label{def:physics:phyx-mini-0062:phyxminiproblems-problemphyxmini0062-mirrorshape}
  \lean{PhyXMiniProblems.ProblemPhyXMini0062.MirrorShape}
  The polygonal reflecting element identified in the diagram
\end{definition}
\begin{proof}
  This is the declared model definition of Mirror Shape.
\end{proof}
\begin{definition}[Mirror Rotation]
  \label{def:physics:phyx-mini-0062:phyxminiproblems-problemphyxmini0062-mirrorrotation}
  \lean{PhyXMiniProblems.ProblemPhyXMini0062.MirrorRotation}
  The axis about which the mirror rotates in the plane of the figure
\end{definition}
\begin{proof}
  This is the declared model definition of Mirror Rotation.
\end{proof}
\begin{definition}[Laser Aim]
  \label{def:physics:phyx-mini-0062:phyxminiproblems-problemphyxmini0062-laseraim}
  \lean{PhyXMiniProblems.ProblemPhyXMini0062.LaserAim}
  The target point of the incident laser ray
\end{definition}
\begin{proof}
  This is the declared model definition of Laser Aim.
\end{proof}
\begin{definition}[Wall Placement]
  \label{def:physics:phyx-mini-0062:phyxminiproblems-problemphyxmini0062-wallplacement}
  \lean{PhyXMiniProblems.ProblemPhyXMini0062.WallPlacement}
  The wall's qualitative placement relative to the laser and mirror
\end{definition}
\begin{proof}
  This is the declared model definition of Wall Placement.
\end{proof}
\begin{definition}[Sweep Endpoint]
  \label{def:physics:phyx-mini-0062:phyxminiproblems-problemphyxmini0062-sweependpoint}
  \lean{PhyXMiniProblems.ProblemPhyXMini0062.SweepEndpoint}
  The two limiting rays produced when the illuminated mirror face changes
\end{definition}
\begin{proof}
  This is the declared model definition of Sweep Endpoint.
\end{proof}
\begin{definition}[Rotating Hexagonal Mirror Setup]
  \label{def:physics:phyx-mini-0062:phyxminiproblems-problemphyxmini0062-rotatinghexagonalmirrorsetup}
  \lean{PhyXMiniProblems.ProblemPhyXMini0062.RotatingHexagonalMirrorSetup}
  \uses{def:physics:phyx-mini-0062:phyxminiproblems-problemphyxmini0062-lengthquantity, def:physics:phyx-mini-0062:phyxminiproblems-problemphyxmini0062-mirrorshape, def:physics:phyx-mini-0062:phyxminiproblems-problemphyxmini0062-mirrorrotation, def:physics:phyx-mini-0062:phyxminiproblems-problemphyxmini0062-laseraim, def:physics:phyx-mini-0062:phyxminiproblems-problemphyxmini0062-wallplacement, def:physics:phyx-mini-0062:phyxminiproblems-problemphyxmini0062-sweependpoint}
  The physical quantities and figure labels for the rotating-mirror setup. The 40 cm arrow in the image runs between opposite vertices of the regular hexagon. mirrorVertexRadius is retained separately because the limiting ray is reflected at the right-hand transition vertex. Wall-hit heights are signed coordinates relative to the laser axis, while streakLengthOnWall is the positive distance between those coordinates
\end{definition}
\begin{proof}
  This is the declared model definition of Rotating Hexagonal Mirror Setup.
\end{proof}
\begin{definition}[Has Depicted Optical Layout]
  \label{def:physics:phyx-mini-0062:phyxminiproblems-problemphyxmini0062-hasdepictedopticallayout}
  \lean{PhyXMiniProblems.ProblemPhyXMini0062.HasDepictedOpticalLayout}
  \uses{def:physics:phyx-mini-0062:phyxminiproblems-problemphyxmini0062-rotatinghexagonalmirrorsetup}
  The categorical arrangement shown by the figure
\end{definition}
\begin{proof}
  This is the declared model definition of Has Depicted Optical Layout.
\end{proof}
\begin{definition}[Matches Figure Readouts]
  \label{def:physics:phyx-mini-0062:phyxminiproblems-problemphyxmini0062-matchesfigurereadouts}
  \lean{PhyXMiniProblems.ProblemPhyXMini0062.MatchesFigureReadouts}
  \uses{def:physics:phyx-mini-0062:phyxminiproblems-problemphyxmini0062-lengthinmeters, def:physics:phyx-mini-0062:phyxminiproblems-problemphyxmini0062-lengthincentimeters, def:physics:phyx-mini-0062:phyxminiproblems-problemphyxmini0062-rotatinghexagonalmirrorsetup}
  The numerical labels read directly from the image: the mirror's vertex-to-vertex width is 40 cm, the center-to-laser distance is 50 cm, and the center-to-wall distance is 2.0 m. No streak length is included
\end{definition}
\begin{proof}
  This is the declared model definition of Matches Figure Readouts.
\end{proof}
\begin{definition}[Has Physical Dimensions]
  \label{def:physics:phyx-mini-0062:phyxminiproblems-problemphyxmini0062-hasphysicaldimensions}
  \lean{PhyXMiniProblems.ProblemPhyXMini0062.HasPhysicalDimensions}
  \uses{def:physics:phyx-mini-0062:phyxminiproblems-problemphyxmini0062-lengthinmeters, def:physics:phyx-mini-0062:phyxminiproblems-problemphyxmini0062-rotatinghexagonalmirrorsetup}
  Positivity and axial ordering of the depicted apparatus. In particular, the laser is outside the mirror and the wall is behind the laser. These physical conditions do not assign a numerical value to the requested streak
\end{definition}
\begin{proof}
  This is the declared model definition of Has Physical Dimensions.
\end{proof}
\begin{definition}[Satisfies Regular Hexagon Transition Geometry]
  \label{def:physics:phyx-mini-0062:phyxminiproblems-problemphyxmini0062-satisfiesregularhexagontransitiongeometry}
  \lean{PhyXMiniProblems.ProblemPhyXMini0062.SatisfiesRegularHexagonTransitionGeometry}
  \uses{def:physics:phyx-mini-0062:phyxminiproblems-problemphyxmini0062-lengthreadout, def:physics:phyx-mini-0062:phyxminiproblems-problemphyxmini0062-rotatinghexagonalmirrorsetup}
  Regular-hexagon transition geometry. Opposite vertices are two circumradii apart. When the incident center ray changes from one face to the next, the outward face normal has reached -π/6 or π/6 from the wall axis
\end{definition}
\begin{proof}
  This is the declared model definition of Satisfies Regular Hexagon Transition Geometry.
\end{proof}
\begin{definition}[Satisfies Law Of Specular Reflection]
  \label{def:physics:phyx-mini-0062:phyxminiproblems-problemphyxmini0062-satisfieslawofspecularreflection}
  \lean{PhyXMiniProblems.ProblemPhyXMini0062.SatisfiesLawOfSpecularReflection}
  \uses{def:physics:phyx-mini-0062:phyxminiproblems-problemphyxmini0062-sweependpoint, def:physics:phyx-mini-0062:phyxminiproblems-problemphyxmini0062-rotatinghexagonalmirrorsetup}
  Specular reflection for the fixed leftward incident ray. With angles measured from the rightward wall axis, reflecting across a face whose normal has angle α produces an outgoing ray at angle 2α
\end{definition}
\begin{proof}
  This is the declared model definition of Satisfies Law Of Specular Reflection.
\end{proof}
\begin{definition}[Satisfies Wall Projection Geometry]
  \label{def:physics:phyx-mini-0062:phyxminiproblems-problemphyxmini0062-satisfieswallprojectiongeometry}
  \lean{PhyXMiniProblems.ProblemPhyXMini0062.SatisfiesWallProjectionGeometry}
  \uses{def:physics:phyx-mini-0062:phyxminiproblems-problemphyxmini0062-lengthreadout, def:physics:phyx-mini-0062:phyxminiproblems-problemphyxmini0062-sweependpoint, def:physics:phyx-mini-0062:phyxminiproblems-problemphyxmini0062-rotatinghexagonalmirrorsetup}
  Planar ray geometry at the vertical wall. At either transition the reflection point is the right-hand mirror vertex, so the horizontal run to the wall is the center-to-wall distance minus the vertex radius. The streak length is the difference between the upper and lower wall-hit coordinates. Both relations are stated in every length unit to keep the model unit-covariant
\end{definition}
\begin{proof}
  This is the declared model definition of Satisfies Wall Projection Geometry.
\end{proof}
\begin{definition}[Answer Choice]
  \label{def:physics:phyx-mini-0062:phyxminiproblems-problemphyxmini0062-answerchoice}
  \lean{PhyXMiniProblems.ProblemPhyXMini0062.AnswerChoice}
  Labels of the four displayed multiple-choice answers
\end{definition}
\begin{proof}
  This is the declared model definition of Answer Choice.
\end{proof}
\begin{definition}[answer Streak Length In Meters]
  \label{def:physics:phyx-mini-0062:phyxminiproblems-problemphyxmini0062-answerstreaklengthinmeters}
  \lean{PhyXMiniProblems.ProblemPhyXMini0062.answerStreakLengthInMeters}
  \uses{def:physics:phyx-mini-0062:phyxminiproblems-problemphyxmini0062-answerchoice}
  The streak length in meters printed beside each answer choice
\end{definition}
\begin{proof}
  This is the declared model definition of answer Streak Length In Meters.
\end{proof}
\begin{definition}[Is Nearest Displayed Streak Length]
  \label{def:physics:phyx-mini-0062:phyxminiproblems-problemphyxmini0062-isnearestdisplayedstreaklength}
  \lean{PhyXMiniProblems.ProblemPhyXMini0062.IsNearestDisplayedStreakLength}
  \uses{def:physics:phyx-mini-0062:phyxminiproblems-problemphyxmini0062-lengthquantity, def:physics:phyx-mini-0062:phyxminiproblems-problemphyxmini0062-lengthinmeters, def:physics:phyx-mini-0062:phyxminiproblems-problemphyxmini0062-answerchoice, def:physics:phyx-mini-0062:phyxminiproblems-problemphyxmini0062-answerstreaklengthinmeters}
  A choice is selected when its printed length is at least as close to the exact physical streak length as every other displayed value. This avoids treating the rounded answer 6.1 m as an exact governing-law hypothesis
\end{definition}
\begin{proof}
  This is the declared model definition of Is Nearest Displayed Streak Length.
\end{proof}
\begin{lemma}[reflected Sweep Endpoint Angles]
  \label{lem:physics:phyx-mini-0062:phyxminiproblems-problemphyxmini0062-reflectedsweependpointangles}
  \lean{PhyXMiniProblems.ProblemPhyXMini0062.reflectedSweepEndpointAngles}
  \uses{def:physics:phyx-mini-0062:phyxminiproblems-problemphyxmini0062-rotatinghexagonalmirrorsetup, def:physics:phyx-mini-0062:phyxminiproblems-problemphyxmini0062-satisfiesregularhexagontransitiongeometry, def:physics:phyx-mini-0062:phyxminiproblems-problemphyxmini0062-satisfieslawofspecularreflection}
  Regular-hexagon transition angles together with specular reflection give a reflected sweep from -π/3 to π/3. This is a derived optical consequence, not a figure readout
\end{lemma}
\begin{proof}
  The conclusion follows from the governing relations and figure readouts named in the statement of reflected Sweep Endpoint Angles.
\end{proof}
\begin{lemma}[streak Length eq hexagon Sweep Formula]
  \label{lem:physics:phyx-mini-0062:phyxminiproblems-problemphyxmini0062-streaklength-eq-hexagonsweepformula}
  \lean{PhyXMiniProblems.ProblemPhyXMini0062.streakLength_eq_hexagonSweepFormula}
  \uses{def:physics:phyx-mini-0062:phyxminiproblems-problemphyxmini0062-lengthinmeters, def:physics:phyx-mini-0062:phyxminiproblems-problemphyxmini0062-rotatinghexagonalmirrorsetup, def:physics:phyx-mini-0062:phyxminiproblems-problemphyxmini0062-satisfiesregularhexagontransitiongeometry, def:physics:phyx-mini-0062:phyxminiproblems-problemphyxmini0062-satisfieslawofspecularreflection, def:physics:phyx-mini-0062:phyxminiproblems-problemphyxmini0062-satisfieswallprojectiongeometry}
  The exact, unrounded streak has twice the upper limiting-ray height. The horizontal run is the wall distance measured from the right-hand transition vertex rather than from the mirror center
\end{lemma}
\begin{proof}
  The conclusion follows from the governing relations and figure readouts named in the statement of streak Length eq hexagon Sweep Formula.
\end{proof}
% --- Archon declaration topology end ---
% --- Archon physics formalization source end ---
